\section{引言}
\label{sec:intro}

\textbf{农业需求：采收环节亟需能“听懂要求、记住进度”的机器人。}采收是果蔬生产中劳动强度最大、季节性最强的环节，其用工高度集中在短短数周的成熟期内。随着农村劳动力持续老龄化和季节性用工日益短缺，“招工难、采收贵、错过最佳采收期”已成为制约设施园艺和果园生产的突出问题，选择性采摘机器人因此长期被视为农业工程的重点方向\cite{bac2014harvesting,kootstra2021selective,zhou2022intelligent}。然而，经过三十余年发展，采摘机器人仍难以大规模进入生产：Bac 等对 50 个采摘机器人项目的综述显示，平均采摘成功率约为 66\%，单果作业周期约 33~s\cite{bac2014harvesting}。更深层的问题在于，现有系统多为“检测—定位—规划—抓取”的模块化流水线\cite{arad2020sweetpepper,williams2019kiwifruit,silwal2017apple}，每一套都针对特定作物、特定决策规则和特定场景布置专门设计，换一种作物、换一种作业要求（例如只摘全红果、按订单摘 20 颗、先摘外侧果）往往意味着重新开发。农业生产真正需要的，是一种能理解作业要求、适应多变场景、在长时作业中保持决策一致的通用采摘智能。

\begin{figure}[t]
\centering
\resizebox{\linewidth}{!}{%
\begin{tikzpicture}[
  font=\footnotesize, >=Stealth,
  inp/.style={draw=black!55, rounded corners=2pt, fill=black!4, align=center, text width=1.9cm, minimum height=0.62cm},
  base/.style={draw=black!45, rounded corners=2pt, fill=black!10, align=center, text width=1.95cm, minimum height=0.62cm},
  oursA/.style={draw=orange!80!black, very thick, rounded corners=3pt, fill=orange!14, align=center, text width=2.7cm},
  oursB/.style={draw=blue!70!black, very thick, rounded corners=3pt, fill=blue!10, align=center, text width=2.7cm},
  oursBs/.style={draw=blue!70!black, very thick, rounded corners=2pt, fill=blue!10, align=center, text width=1.95cm, minimum height=0.62cm},
  oursC/.style={draw=green!50!black, very thick, dashed, rounded corners=3pt, fill=green!8, align=center, text width=1.9cm},
  outn/.style={draw=black!55, rounded corners=2pt, fill=white, align=center, text width=1.8cm, minimum height=0.62cm},
  lab/.style={font=\scriptsize, fill=white, inner sep=1pt}]

\node[inp] (lang) at (1.2,6.7) {语言指令 $\ell$};
\node[inp] (img)  at (1.2,5.85) {图像 $I^{f}_t,\,I^{w}_t$};
\node[inp] (st)   at (1.2,5.0) {本体状态 $s_t$};
\node[oursC, minimum height=1.4cm] (c3) at (1.2,3.1) {\textbf{\Ciii{} 事件鲁棒训练}\\\textbf{与免标注}\\\textbf{真机适配}};
\node[inp, minimum height=0.9cm] (buf) at (1.2,1.3) {历史帧缓存\\$I_{0:t}$};

\node[oursA, minimum height=1.3cm] (pedet) at (5.0,3.9) {\textbf{\Ci{} 采摘事件检测器}\\\textbf{PEDet}\\{\scriptsize 腕部视觉 + 本体感知}\\{\scriptsize 因果 TCN，区分成败}};
\node[oursB, minimum height=1.8cm] (hem) at (5.0,1.3) {\textbf{\Cii{} 事件中心分层记忆}\\\textbf{HEM}\\{\scriptsize 场景锚 / 工作事件 /}\\{\scriptsize 情景摘要 / 时间上下文}};

\begin{scope}[on background layer]
\node[draw=black!40, dashed, rounded corners=4pt, fill=black!2, minimum width=8.6cm, minimum height=7.0cm] (vla) at (11.85,3.75) {};
\end{scope}
\node[font=\scriptsize, text=black!60] at (11.85,0.5) {基础 VLA 结构（灰色模块沿用，不作为本文贡献）};
\node[base] (tok)  at (8.9,6.7) {分词器};
\node[base, minimum height=0.75cm] (enc1) at (8.9,5.85) {视觉编码器\\{\scriptsize SigLIP（冻结）}};
\node[base, minimum height=1.6cm] (vlm) at (12.0,6.275) {VLM 骨干\\{\scriptsize Gemma-2B}};
\node[base, minimum height=1.2cm] (act) at (14.85,3.7) {动作专家\\{\scriptsize 流匹配}};
\node[base, minimum height=0.9cm] (mexp) at (14.85,1.3) {记忆分支};
\node[base, minimum height=0.9cm] (enc2) at (8.9,1.3) {视觉编码器\\{\scriptsize 共享，16 token/帧}};
\node[oursBs, minimum height=0.9cm] (emb) at (12.0,1.3) {\textbf{\Cii{} 类型-时间}\\\textbf{嵌入}};

\node[outn, minimum height=0.9cm] (chunk) at (17.6,3.7) {动作块 $A_t$\\{\scriptsize $20\times8$，执行 16 步}};
\node[outn] (guard) at (17.6,2.5) {安全包络};
\node[outn] (arm)   at (17.6,1.3) {机械臂 20~Hz};

\draw[->] (lang) -- (tok);
\draw[->] (img) -- (enc1);
\draw[->] (tok) -- (vlm.west |- tok);
\draw[->] (enc1) -- (vlm.west |- enc1);
\draw[->] (5.8,5.85) -- node[lab, right]{腕部} (5.8,4.55);
\draw[->] (st.east) -- (4.2,5.0) -- (4.2,4.55);
\draw[->, very thick, orange!80!black] (pedet) -- node[lab, right]{事件表 $\mathcal E_t$} (hem);
\draw[->] (buf) -- (hem.west |- buf);
\draw[->, very thick, blue!70!black] (hem.east |- enc2) -- node[lab, above]{$B{=}16$ 幅} (enc2);
\draw[->] (enc2) -- (emb);
\draw[->, very thick, blue!70!black] (emb) -- (mexp);
\draw[->] (vlm.east) -| (act.north);
\draw[->] (mexp) -- (act);
\draw[->] (act) -- (chunk);
\draw[->] (chunk) -- (guard);
\draw[->] (guard) -- (arm);
\draw[->, dashed, green!45!black, thick] (c3.east |- 0,3.6) -- node[lab, above]{伪标签} (pedet.west |- 0,3.6);
\draw[->, dashed, green!45!black, thick] ([yshift=0.25cm]c3.south east) -- (3.05,2.65) |- node[lab, pos=0.75, above]{扰动事件} (hem.west |- 0,1.9);

\node[font=\scriptsize, anchor=north] at (9.4,-0.05) {\sq{orange!80!black}\ \Ci{} 何时发生了采摘\qquad \sq{blue!70!black}\ \Cii{} 记住什么、如何编码\qquad \sq{green!45!black}\ \Ciii{} 如何对检测误差鲁棒、如何零标注上真机};
\end{tikzpicture}}
\caption{HarvestMem 总体框架。灰色为通用 VLA 的沿用结构：冻结的 SigLIP 编码当前图像，VLM 骨干融合指令（本体状态以离散化 token 一并输入，图中省略），流匹配动作专家输出 20 步动作块。彩色为本文创新：\Ci{} \pedet{} 从腕部图像和本体状态因果地检测成功采摘；\Cii{} \hem{} 依据事件表把 16 幅记忆的预算分配给场景锚、工作事件、情景摘要拼图帧和时间上下文，并在共享视觉编码器之后叠加类型-时间嵌入；\Ciii{} 训练时按检测器误差模型扰动事件，部署时由作业结果事后确认生成伪标签、持续更新 \pedet。记忆图像数和 token 数与基础结构完全相同。}
\label{fig:overview}
\end{figure}


\textbf{通用 VLA 带来了新范式。}视觉-语言-动作（vision-language-action，VLA）模型把在互联网规模图文数据上预训练的视觉-语言模型与大规模机器人示范结合，直接把图像、本体状态和自然语言指令映射为动作\cite{brohan2023rt1,brohan2023rt2,octo2024,kim2024openvla,black2024pi0,pi05_2025}。在这一范式下，“摘哪些果、摘几颗、放到哪里”只需写进指令，感知、决策与控制在一个网络中端到端完成，同一模型可以跨任务、跨场景迁移。这为构建“一个模型、多种作业”的通用采摘机器人提供了可能。

\textbf{但通用 VLA 并不能直接胜任农业作业，农业场景需要“特需 VLA”。}现有 VLA 主要在实验室桌面任务上发展起来，这类任务时间短（通常数十秒）、场景静态、目标始终可见。农业采摘则截然不同，具有四个通用 VLA 未曾面对的特点：
\begin{enumerate}
\item \textbf{动态遮挡。}风、机械臂自身的运动和叶片回弹会让片刻之前清晰可见的成熟果实被完全遮住\cite{hemming2014detectability}。在本文的仿真基准中，被叶片遮住的成熟果在仍挂在树上的帧里，中位数有 \pct{occl_invis} 的时间完全不可见。
\item \textbf{作业进度不可见。}集果篮通常位于相机视野之外，机器人无法通过“再看一眼”确认已经摘了几颗；按配额采收、分级抽样和疏果都要求准确记住自己做过什么。
\item \textbf{失败与重试频繁。}果梗韧性不一、果实易滑脱，抓空和重试在真实作业中极为常见，机器人必须记住哪些采摘\emph{真正完成了}，而不是哪些动作\emph{尝试过}。
\item \textbf{长时、重复、算力受限。}一株植物、一行植株的作业动辄数百至数千个控制步，而田间平台的机载算力有限，无法承受随时长增长的计算开销。
\end{enumerate}
这四点共同指向同一项能力缺口：\textbf{在固定计算预算下，记住长时作业中稀疏而关键的交互事件}。遗忘的农艺代价是直接的：忘掉一颗被遮住的成熟果，就是留在树上过熟的产量损失；忘掉一次已完成的采摘，就会超量采收；把失败尝试当成功，则会漏采。

\textbf{现有 VLA 的记忆机制与农业需求存在结构性错配。}具备记忆的 VLA 只能保留有限帧数的历史画面，因为每一帧都会增加视觉 token，注意力计算量随 token 数平方增长\cite{vaswani2017attention}；在这一预算内，历史帧通常按时间均匀采样或只保留最近窗口。我们发现，这一设计对采摘作业存在根本缺陷：定义一次采摘的画面，即夹爪包住果实和果实脱离果柄，只持续约 \V{pick_frames} 帧，均匀采样得到的帧中 TCP 距待摘果 3~cm 以内的不足 \pct{near_frac}。第~\ref{sec:coverage} 节将证明，均匀记忆对已完成采摘的覆盖率随作业时长按 $1/t$ 衰减。实测中，16 帧均匀记忆漏掉了 \pct{cov_u16_miss} 的已完成采摘，只有 \pct{cnt_u16} 的决策时刻能从记忆中数清已摘数量；作业超过 600 帧后，连最近一次采摘也只在 \pct{last_u_late} 的决策中被保留。换言之，无论已经摘了多少颗，均匀记忆平均只“记得”约 \V{cap_u} 次采摘。

\textbf{直观的补救办法都不能解决问题。}扩大记忆预算只是治标：要维持覆盖率，均匀记忆的帧数必须随作业时长线性增长，计算开销按平方增长。提升单帧感知能力看不见被遮挡的果实，也会把夹爪中正在搬运的果实误当作新目标；在我们的实验中，逐帧检测只在 \pct{perframe_cnt} 的回合中得到正确的采摘计数。直接为最近几次采摘预留关键帧可以缓解问题，但每次采摘占用两帧，16 帧预算最多只能覆盖 \V{cap_flat_max} 次采摘且不再留有任何时间上下文，远不能满足一行植株数十次采摘的需求。而无论哪种基于事件的方案，都面临一个现实问题：在真机上，事件从哪里来？仿真中的果实位置真值不可得，手工感知链路对光照敏感（光照变化下分割 F1 低于 \V{hsv_light_f1}），而在果实上加装传感器在真实生产中并不可行。

\textbf{本文思路：让 VLA 像人一样按事件组织记忆。}认知科学研究表明，人会把连续经历切分为离散事件，并优先编码事件边界\cite{zacks2007event}；情景记忆把“发生了什么”与“何时、何地”绑定在一起\cite{tulving1972episodic,clayton1998episodic}。熟练的采摘工人不会记住每一秒看到的画面，但会清楚记得“刚才那颗摘下来了”“这棵树已经摘了五颗”。受此启发，本文提出面向农业采摘的事件中心记忆 VLA 模型 \textbf{HarvestMem}（图~\ref{fig:overview}）。它在通用 VLA 的结构之上，只替换“记什么”这一环节：由学习型检测器从机器人自身的传感器流中识别采摘事件，再由分层记忆把有限的记忆预算优先分配给事件，并以压缩形式长期保留远期事件。据我们所知，这是首个以物理交互事件组织 VLA 记忆、并面向农业长时采摘作业系统验证的工作。

本文的主要贡献如下：
\begin{enumerate}
\item \textbf{揭示通用 VLA 记忆与农业作业的结构性错配。}把农业 VLA 的记忆问题形式化为固定预算下的事件覆盖问题，证明均匀记忆的事件覆盖率 $c_U(t)\approx\min(1,(B-1)|S|/t)$ 随作业时长反比衰减（命题~\ref{prop:uniform}），理论预测与实测误差不超过 \V{theory_err}。
\item \textbf{创新一：物理监督的学习型采摘事件检测器 \pedet。}以“双指接触 + 果柄脱离”的物理定义作为监督信号，在仿真中零人工标注地自动生成逐帧标签；融合腕部图像与本体感知的因果时序网络能区分成功采摘与失败尝试，不依赖任何特权信息，对光照变化鲁棒（光照扰动下 F1 \V{det_pe_light}，手工链路为 \V{det_fdm_light}）。
\item \textbf{创新二：事件中心分层记忆 \hem。}把固定的 $B$ 帧记忆组织为场景锚、工作事件、情景摘要和时间上下文四层，提出“事件拼图帧”把多次远期采摘压缩进一个记忆位置，并以零初始化的类型-时间嵌入告诉策略每帧记忆“是什么、多久以前”。在记忆帧数和推理开销完全不变的条件下，16 帧预算可保留的采摘次数从约 \V{cap_u} 次（均匀）和约 \V{cap_k8} 次（平坦关键帧）提升到 \V{cap_hem} 次（命题~\ref{prop:capacity}）。
\item \textbf{创新三：事件鲁棒训练与免标注真机适配。}训练时按检测器的误差模型扰动事件，使策略对漏检、误检和时间偏差具有鲁棒性；部署时提出“事后确认自训练”，仅凭夹爪开度与腕部相机观察到的作业结果（果实是否放入篮中）自动生成伪标签，无需人工标注，也无需在果实或植株上加装任何传感器，使检测器可在作业中持续自适应。
\item \textbf{经过捷径审计的农业记忆基准与系统验证。}构建 \bench{} 基准，提出用信息受限的脚本策略审计任务捷径，设计分别依赖事件记忆、场景记忆和长时事件记忆的三类任务，并在封存协议下完成仿真配对评测与真实番茄采摘实验。
\end{enumerate}

本文其余部分安排如下：第~\ref{sec:related} 节综述相关工作；第~\ref{sec:problem} 节形式化问题并分析均匀记忆的覆盖率；第~\ref{sec:method} 节介绍 HarvestMem 的数据链路与三项创新；第~\ref{sec:benchmark} 节介绍基准及其有效性审计；第~\ref{sec:design} 节给出仿真与真机实验设计；第~\ref{sec:results} 节报告结果；第~\ref{sec:discussion} 节讨论意义与局限并总结全文。
